\documentclass[reqno, 12pt]{amsart}

\usepackage{amssymb}
\usepackage{amsthm}
\usepackage{amsmath}
\usepackage{amsxtra}
\usepackage{mathrsfs}
\usepackage{comment}
\usepackage{graphicx}
\usepackage{placeins}
\usepackage{multicol}
\usepackage{bbm}
\usepackage{stmaryrd}
%\usepackage{enumerate}
\usepackage[inline, shortlabels]{enumitem}
\usepackage{mathdots}
\usepackage{colonequals}
\usepackage{hyperref}


\usepackage{calc}
\usepackage{tikz}
\usetikzlibrary{arrows,calc,automata,shadows,backgrounds,positioning,intersections,fadings,decorations.pathreplacing,shapes,matrix}
\usepackage{tikz-cd}

\usepackage{fullpage}

\newtheorem{thm}{Theorem}[section]
\newtheorem{lem}[thm]{Lemma}
\newtheorem{cor}[thm]{Corollary}
\newtheorem{conj}[thm]{Conjecture}
\newtheorem{prop}[thm]{Proposition}
\theoremstyle{definition}  
\newtheorem{defn} [thm] {Definition} 
\newtheorem{claim}[thm]{Claim}
\newtheorem{example} [thm] {Example}
\newtheorem{rem} [thm] {Remark}

\newenvironment{problab}[1]
{\noindent\textbf{Problem #1}.}
{\vskip 6pt}

\newenvironment{exolab}[1]
{\noindent\textbf{Exercise #1}.}
{\vskip 6pt}

\theoremstyle{remark}
\newtheorem*{soln}{Solution}

\newcommand{\C}{\mathbb C}
\renewcommand{\H}{\mathbb H}
\newcommand{\D}{\mathbb D}
\newcommand{\F}{\mathbb F}
\newcommand{\Q}{\mathbb Q}
\newcommand{\R}{\mathbb R}
\newcommand{\Z}{\mathbb Z}
\newcommand{\T}{\mathbb T}
\renewcommand{\P}{\mathcal P}
\renewcommand{\O}{\mathcal O}
\newcommand{\m}{\mathfrak m}
\newcommand{\A}{\mathbb A}
\renewcommand{\SS}{\mathbb S}
\renewcommand{\L}{\mathcal L}

\newcommand{\B}{\mathcal B}
\newcommand{\E}{\mathcal E}

\newcommand{\SSpan}{\text{span}}

\DeclareMathOperator{\lcm}{lcm}
\DeclareMathOperator{\img}{img}
\DeclareMathOperator{\Ann}{Ann}
\DeclareMathOperator{\Tor}{Tor}
\DeclareMathOperator{\tr}{tr}
\DeclareMathOperator{\Hom}{Hom}
\DeclareMathOperator{\End}{End}
\DeclareMathOperator{\Aut}{Aut}
\DeclareMathOperator{\Gal}{Gal}
\DeclareMathOperator{\sgn}{sgn}
\DeclareMathOperator{\ord}{ord}
\DeclareMathOperator{\ad}{ad}
\DeclareMathOperator{\coker}{coker}
\DeclareMathOperator{\rank}{rank}

\DeclareMathOperator{\id}{id}

\usepackage{mathpazo}

\everymath{\displaystyle}

\tikzset{commutative diagrams/.cd, arrow style = tikz, diagrams = {>=latex}}

\newenvironment{amatrix}[1]{%
  \left(\begin{array}{@{}*{#1}{c}|c@{}}
}{%
  \end{array}\right)
}

\usepackage{abstract}
\renewcommand{\abstractname}{}    % clear the title
\renewcommand{\absnamepos}{empty} % originally center
%\setlength{\abstitleskip}

\begin{document}
%\renewcommand{\abstractname}{\vspace{-\baselineskip}}

\title{18.700 Problem Set 5}
%\author{Évariste Galois}
\dedicatory{Due Wednesday, October 23 at 11:59 pm on \href{https://canvas.mit.edu/courses/27315}{Canvas}}
%\thanks{}
%\contrib{}
%\address{}

\maketitle

\begin{abstract}
  \noindent Collaborated with:\\
  Sources used: 
\end{abstract}
\vskip 0.5cm

\begin{problab}{1} (3D \#18) (8 points)
  Show that $V$ and $\L(\F, V)$ are isomorphic $\F$-vector spaces.
\end{problab}

%\begin{soln}
%\end{soln}

\begin{problab}{2} (3D \#6) (10 points)
  Suppose that $W$ is finite-dimensional and $S,T \in \L(V,W)$. Prove that $\ker(S) = \ker(T)$ if and only if there exists an invertible $P \in \L(W)$ such that $S = PT$.
\end{problab}

\begin{problab}{3} (5A \#20) (8 points)
  Define the left shift operator $L \in \L(\F^\infty)$ by
  \begin{equation*}
    L(z_1, z_2, z_3, \ldots) = (z_2, z_3, \ldots) \, .
  \end{equation*}
  \begin{enumerate}[(a)]
    \item 
      Show that every element of $\F$ is an eigenvalue of $L$.
    \item
      Find all eigenvectors of $L$.
  \end{enumerate}
\end{problab}

\begin{problab}{4} (5A \#21) (8 points)
  Suppose $T \in \L(V)$ is invertible.
  \begin{enumerate}[(a)]
    \item 
      Suppose $\lambda \in \F$ with $\lambda \neq 0$. Prove that $\lambda$ is an eigenvalue of $T$ if and only if $1/\lambda$ is an eigenvalue of $T^{-1}$.
    \item
      Prove that $T$ and $T^{-1}$ have the same eigenvectors.
  \end{enumerate}
\end{problab}

\begin{problab}{5} (5B \#11) (10 points)
  Suppose $V$ is a $2$-dimensional vector space, $T \in \L(V)$, and the matrix of $T$ with respect to some basis of $V$ is
  $
  \begin{pmatrix}
    a & b\\
    c & d
  \end{pmatrix}.
  $
  \begin{enumerate}[(a)]
    \item 
      Show that $T^2 - (a+d)T + (ad-bc)I = 0$.
    \item
      Show that the minimal polynomial of $T$ is
      \begin{equation*}
        \begin{cases}
          z-a & \text{if $b=c=0$ and $a=d$,}\\
          z^2 - (a+d)z + (ad-bc) = 0 & \text{otherwise.}

        \end{cases}
      \end{equation*}
  \end{enumerate}
\end{problab}

\begin{problab}{6} (5A \#30) (5 points)
  Suppose $T \in \L(V)$ and
  \begin{equation*}
    (T-2I)(T-3I)(T-4I) = 0 \, .
  \end{equation*}
  Suppose $\lambda$ is an eigenvalue of $T$. Prove that $\lambda = 2$ or $\lambda = 3$ or $\lambda = 4$.
\end{problab}

\end{document} 	
